\section{Basic (Python)}

\subsection{Small Hacks}

\begin{center}
\begin{tabular}{lll}
\hline
Expression & C++ result & Python result \\
\hline
\texttt{-7 // 2}  & $     -3$ & $     -4$ \\
\texttt{-7 \% 2}  & $     -1$ & $\;\;\:1$  \\
\texttt{7 // -2}  & $     -3$ & $     -4$ \\
\texttt{7 \% -2}  & $\;\;\:1$ & $     -1$ \\
\texttt{-7 // -2} & $\;\;\:3$ & $\;\;\:3$  \\
\texttt{-7 \% -2} & $     -1$ & $     -1$ \\
\hline
\end{tabular}
\end{center}

\pycode{basic/hacks/io.py}
\pycode{basic/hacks/recursion.py}
\pycode{basic/hacks/itertools.py}

\subsection{Find First Ok (Binary Search)}
Binary search on a monotone predicate (False\dots False True\dots True). Searches \verb|[lb, rb)|, returns the first \verb|x| with \verb|ok(x) = True|, or \verb|rb| if there is none.
\tc{\log{(rb - lb)}}
\pycode{basic/find_first_ok.py}

\subsection{Offline Segments Sum}
Difference array: add \verb|x| to every position of \verb|[l, r)| in \tc{1}per update, then one prefix-sum pass gives the final array. 
\tc{n + q}
\pycode{basic/offline_segments_sum.py}

\subsection{Compress Values}
Replaces each value by its rank among distinct sorted values; also returns the sorted distinct values to map back. 
\tc{n \log n}
\pycode{basic/compress_values.py}

\subsection{Heap Example}
heapq is a min-heap on a list: \verb|heappush| / \verb|heappop| are \tc{\log n}.
\pycode{basic/heap_example.py}

\subsection{Previous Smaller}
For each \verb|i|, the index of the nearest element to the left that is strictly smaller, or \verb|-1| (monotonic stack). 
\tc{n}
\pycode{basic/previous_smaller.py}

\subsection{Window Min}
Minimum of every window of length \verb|k| via a monotonic deque of indices.
\tc{n}
\pycode{basic/window_min.py}

\subsection{DSU}
Disjoint set union with iterative path compression and random linking (\tc{1} in fact). \verb|join| returns \verb|False| if both elements were already in the same set.
\pycode{basic/dsu.py}

\subsection{Fenwick Tree}
Point update and range sum in \tc{\log n}, built in \tc{n}. Indices are 0-based; \verb|query(lb, rb)| is end-inclusive.
\pycode{basic/fentree.py}

\subsection{DFS Iterative}
Depth-first search with an explicit stack of iterators (no recursion limit problems). Returns the vertices in preorder.
\tc{n + m}
\pycode{basic/dfs_iterative.py}

\subsection{BFS}
Shortest paths by edge count from \verb|s|. Returns the distance (-1 if unreachable) and the previous vertex for path restoration.
\tc{n + m}
\pycode{basic/bfs.py}

\subsection{BFS 0-1}
Shortest paths when edge weights are only 0 or 1: 0-edges go to the front of the deque, 1-edges to the back; stale entries are skipped.
\tc{n + m}
\pycode{basic/bfs_01.py}

\subsection{Graph Components}
Labels each vertex with its connected component using iterative DFS; returns the labels and the number of components.
\tc{n + m}
\pycode{basic/graph_components.py}

\subsection{Kruskal (Min Spanning Tree)}
Minimum spanning tree: sort edges by weight, add an edge if it joins two different components (uses \verb|DSU|). Returns \verb|(total, edges)| or \verb|None| if the graph is disconnected.
\tc{n + m}
\pycode{basic/kruskal.py}

\subsection{Knapsack 0-1}
Each item is used at most once; the inner loop over capacity goes downward so an item is not reused.
\tc{n \cdot \text{capacity}}
\pycode{basic/knapsack_01.py}

\subsection{Knapsack Unbounded}
Each item can be used any number of times; the inner loop goes upward.
\tc{n \cdot \text{capacity}}
\pycode{basic/knapsack_unbounded.py}

\subsection{Min Coins}
Minimum number of coins summing to target, or -1 if impossible.
\tc{\text{target} \cdot \text{len(coins)}}
\pycode{basic/min_coins.py}

\subsection{Longest Increasing Subsequence}
Keeps the smallest possible tail for each length and replaces tails via binary search. Strict version with \verb|bisect_left|, non-decreasing with \verb|bisect_right|.
\tc{n \log n}
\pycode{basic/lis_length.py}

\subsection{Longest Common Subsequence}
Length of the LCS with a rolling DP row.
\tc{\text{len(a)} \cdot \text{len(b)}}
\pycode{basic/lcs_length.py}

\subsection{Levenshtein Editing Distance}
Minimum number of insertions, deletions and substitutions to turn \verb|a| into \verb|b|.
\tc{\text{len(a)} \cdot \text{len(b)}}
\pycode{basic/levenshtein.py}

\subsection{Binary Pow}
Computes $a^b \% m$ in \tc{\log b} by repeated squaring (the built-in \verb|pow(a, b, mod)| does the same).
\pycode{basic/binpow.py}

\subsection{Binomials}
Precomputes factorials and inverse factorials in \tc{\text{n\_max}} and returns \verb|choose(n, k) % mod| in \tc{1}. \verb|mod| must be prime and greater than \verb|n_max|.
\pycode{basic/binomial_table.py}